\section{Pinned source bridge and quantifier ledger}\label{app:bridge}

This appendix states exactly what is imported from the Gaussian source
construction and what is proved locally.  It is included to make the claim
boundary auditable without relying on status labels in the research log.

\subsection{Version-pinned inputs}

The source snapshots are Hunter--Milojevi\'c--Sudakov v2
\cite{HMS2025} and Lin--Niu v2 \cite{LinNiu2026}.  Their source hashes,
filenames, and retrieval URLs are recorded in the repository manifest.  The
imported ingredients are:
\begin{enumerate}[label=(S\arabic*)]
 \item the HMS Gaussian construction and Appendix-B reverse induction;
 \item the red perfect-sequence extraction and the original blue bound;
 \item the centered square-moment estimate in the range used in (A.16);
 \item the one-sided truncated-Gaussian cumulant identity used to define
       $G_*(C)$; and
 \item the first-moment conversion from red and blue probability exponents to
       a Ramsey lower bound.
\end{enumerate}
The local argument changes none of these statements.  It adds a deterministic
negative term to the red ledger, checks that extraction retains it, and then
uses the same red--blue comparison.

\subsection{Literal red-side bridge}

For orientation, the preceding source audit establishes a red deficit of the
form
\[
 \mathcal D_C(r)=\frac{c_R(C)}d\sum_{k=1}^{r-1}k(k-1)^2,
 \qquad c_R(C)>0, \tag{B.1}
\]
inside the original Appendix-B recursion.  Since
\[
 \sum_{k=1}^{r-1}k(k-1)^2
 =\frac{r(r-1)(r-2)(3r-5)}{12}, \tag{B.2}
\]
the corresponding rate term is $c_R(C)/(4D^2)$.  This bridge proves that a
same-history, pre-coarsening deficit can be carried through the public
induction and extraction.  The optimized exact-CGF term $G_*(C)$ used in
the present paper is the reviewed refinement of that bridge.

The blue-side construction is not modified here.  For sufficiently large
fixed $C$, its source exponent remains strictly above the red exponent,
because every added red term is $O(1/\log C)=o(1)$ whereas the frozen
red--blue gap is of constant order before the final $C\to\infty$ limit.
Thus the red side remains the bottleneck.

\subsection{Order of quantifiers}

The source $O(\cdot)$ and $o(\cdot)$ terms do not carry published numerical
constants.  The formal theorem therefore has the order
\[
 \exists C_0\ \forall C\ge C_0\ \exists\ell_0(C)\
 \forall\ell\ge\ell_0(C), \tag{B.3}
\]
followed by the limit $\ell\to\infty$.  Only after this fixed-$C$ theorem
is established do we take $C\to\infty$ to obtain the coefficient $1/64$.
No explicit value of $C_0$ or $\ell_0(C)$ is asserted.

The source error constant $K$ is frozen before $C_0$ is selected.  The
new ledger deletes probability mass relative to the source envelope and adds
only the explicitly bounded deterministic error $J_{r,d}$ in (A.23).
Therefore no new, history-dependent source constant is introduced.

\subsection{Claim boundary}

The theorem is source-relative in two senses.  First, it is conditional on
the version-pinned interfaces (S1)--(S5).  Second, comparison with a printed
final constant in another normalization requires translating both estimates
from the same probability inequality.  We have not used such a translation
to make a priority or best-known claim.  What is proved here is the internal
strengthening
\[
 G_*(C)+H_*(C)\quad\longrightarrow\quad
 G_*(C)+\widehat H_*(C),\qquad
 \widehat H_*(C)>H_*(C)
\tag{B.4}
\]
for all sufficiently large fixed $C$, together with
$\widehat H_*(C)=(1+o(1))/(64\log C)$.

The frozen arithmetic checker independently verifies the triangular
recursion, the exact $K_r$ identity, the full-Hessian inequalities at
representative large values of $C$, the strict positivity margins, and the
asymptotic coefficient.  It does not replace the probabilistic induction;
that induction is given in \cref{sec:induction,app:one-step}.
